\section{Tóm tắt nghiên cứu}
\label{sec:abstract}

Bài toán Dự báo Phụ tải Điện Ngắn hạn (\textit{Short-Term Load Forecasting} -- STLF) theo từng giờ có ý nghĩa sống còn đối với công tác điều độ kinh tế, tối ưu hóa công suất phát điện dự phòng (\textit{spinning reserve}) và đảm bảo an ninh vận hành lưới điện khu vực của \textit{Pacific Gas and Electric Company} (PG\&E) tại California. Tuy nhiên, các phương pháp hiện có, bao gồm cả mô hình đối chuẩn hàng đầu của Roy và cộng sự~\cite{roy2025hourly} tại cuộc thi \textit{IISE PG\&E Energy Analytics Challenge 2025}, còn bộc lộ ba hạn chế lớn: thiếu vắng tri thức nhiệt động lực học công trình (\textit{thermodynamic inductive bias}) tại ngưỡng kích hoạt điều hòa nhiệt độ $18^\circ\text{C}$, các biến phân loại nhị phân rời rạc theo tháng gây ra hiện tượng đứt gãy gián đoạn tại ranh giới chuyển mùa, và việc áp dụng một họ mô hình cây đơn lẻ đã bỏ ngỏ tiềm năng đa dạng sai số (\textit{error diversity}) giữa các lớp kiến trúc khác nhau dưới các chế độ vận hành lưới điện. 

Nghiên cứu này đề xuất khung phương pháp \textbf{Context-Aware Dual-Regime Stacking (CAS)} tích hợp kỹ nghệ đặc trưng định hướng vật lý nhằm trả lời ba câu hỏi nghiên cứu về hiệu quả của đặc trưng nhiệt động lực học, tính bù trừ sai số hai chế độ vận hành và năng lực định tuyến ngữ cảnh thời gian thực. Thí nghiệm được thực hiện trên tập dữ liệu chuỗi thời gian thực tế 3 năm (2020--2022) của PG\&E với 26,304 điểm quan sát và so sánh đối chuẩn trực tiếp với mô hình 24-Hourly XGBoost của Roy et al.~\cite{roy2025hourly}. 

Kết quả thực nghiệm trên tập kiểm tra độc lập năm thứ ba (Year 3 / 2022, 8,760 giờ) cho thấy mô hình CAS đề xuất đạt kỷ lục mới với RMSE \textbf{141.81~MW}, MAPE \textbf{4.75\%}, và $R^2$ đạt \textbf{0.9136}, cắt giảm tới \textbf{$-$29.03~MW (tương đương giảm $-$17.0\% RMSE)} so với baseline bài báo gốc với ý nghĩa thống kê vượt trội qua kiểm định Diebold-Mariano ($t > 8.5, p < 10^{-5}$). Phân tích bổ sung chỉ ra sự phân tầng hiệu năng rõ nét: mô hình tuyến tính ElasticNet (L1/L2 regularization) vượt trội hơn XGBoost trong chế độ phụ tải nền ban đêm (00:00--06:00: RMSE $131.33$~MW so với $141.37$~MW), trong khi XGBoost chiếm ưu thế áp đảo trong chế độ sóng nhiệt mùa hè (CDD $> 3^\circ\text{C}$: RMSE $244.16$~MW so với $313.27$~MW); mô hình meta-learner LightGBM nhận biết ngữ cảnh đã học được quy luật định tuyến động để vượt trội hơn cả hai mô hình thành phần trên toàn bộ 8 chế độ vận hành thực tế. 

Trong phạm vi điều độ lưới điện khu vực, kết quả cho thấy việc giảm $29.03$~MW sai số RMSE giúp tiết kiệm ước tính từ \textbf{5.0 đến 7.6 triệu USD mỗi năm} chi phí mua công suất dự phòng quay trên thị trường điện California (CAISO), đồng thời cung cấp giải pháp dự báo tin cậy và minh bạch phục vụ chuyển dịch năng lượng bền vững.
